% Generated draft: requires mathematical and source review.
\documentclass[11pt]{article}
\usepackage[margin=25mm]{geometry}
\usepackage{fontspec}
\setmainfont{Latin Modern Roman}
\usepackage{amsmath,amssymb,mathtools,unicode-math}
\setmathfont{Latin Modern Math}
\usepackage{enumitem,needspace,xurl,hyperref}
\setlength{\parindent}{0pt}
\setlength{\parskip}{5pt}
\setlength{\emergencystretch}{3em}
\newcommand{\N}{\mathbb{N}}
\newcommand{\Z}{\mathbb{Z}}
\newcommand{\Q}{\mathbb{Q}}
\newcommand{\R}{\mathbb{R}}
\newcommand{\C}{\mathbb{C}}
\newcommand{\F}{\mathbb{F}}
\newcommand{\A}{\mathbb{A}}
\newcommand{\PP}{\mathbb{P}}
\newcommand{\E}{\mathbb{E}}
\newcommand{\eps}{\varepsilon}
\newcommand{\dd}{\,\mathrm{d}}
\begin{document}
\section*{UnsolvedMath ID: 800762}
\textbf{Source problem code:} catalogue entry
\subsection{Definitions and mathematical statement}
Let \(K\) range over fields, with no characteristic restriction stated in the supplied record. In the standard curve formulation, \(K\) is called \emph{ample} if every smooth geometrically integral curve \(C\) over \(K\) satisfies
\[
C(K)\ne\varnothing\quad\Longrightarrow\quad \#C(K)=\infty,
\]
where \(C(K)\) denotes its set of \(K\)-rational points.

Write \(A(K)\), \(V(K)\), and \(W(K)\) for the assertions that \(K\) is ample, virtually ample, and weakly ample, respectively. The supplied record does not define \emph{virtually ample} or \emph{weakly ample}. Consequently, \(V\) and \(W\) remain placeholders for the source's intended properties, and the questions below are not yet a self-contained mathematical formulation.

The source states the implications
\[
\forall K,\qquad A(K)\Longrightarrow V(K)\Longrightarrow W(K).
\]
Its Question 5 asks for the relation among these three properties and, specifically, whether
\[
\forall K,\qquad W(K)\Longrightarrow V(K),
\]
and whether
\[
\forall K,\qquad V(K)\Longrightarrow A(K).
\]
The source presents analogous descent properties for Henselian valued fields as motivation. Its assertion that no proof of the second converse is known is part of the supplied historical account, not an independently verified current literature assessment.
\subsection{Short English statement}
Bary-Soroker and Fehm ask whether every weakly ample field is virtually ample and whether every virtually ample field is ample, and more generally how these properties are related. The missing definitions of virtually ample and weakly ample must be recovered before this draft is self-contained.
\subsection{Sources}
\begin{itemize}
\item Local UnsolvedMath registry, permanent ID \(800762\), source problem code \texttt{catalogue entry}, titled \emph{Open Problems in the Theory of Ample Fields --- Question 5}. The registry attributes the problem to Bary-Soroker and Fehm, with proposed year 2011, and labels it \emph{open}; this label is supplied metadata, not independently verified mathematical status.
\item Bary-Soroker and Fehm, \emph{Open Problems in the Theory of Ample Fields} (2011), Question 5. Supplied source URL: \url{https://arxiv.org/abs/1106.1310}. The supplied provenance records access on 2026-10-16 and extraction by \texttt{arxiv\_latex}. These are supplied provenance claims; no source visit or independent verification is asserted here.
\end{itemize}
\end{document}
